% =============================================================================
%  §2  Ch.5 循环单调性与 Kantorovich 对偶
%  来源：Villani, Optimal Transport: Old and New, Chapter 5 (pp. 51–92)
% =============================================================================
\TLsection{2 \ 循环单调性与 Kantorovich 对偶}

% ── 本节要回答什么问题 ──────────────────────────────────────────────────────────
\begin{frame}{本节：最优传输的两个核心工具}
  第 1 节证明了最优耦合\emph{存在}，但对其\textbf{结构}只字未提。本节引入两个互补视角：

  \begin{itemize}
    \item \textbf{循环单调性}（几何视角）：最优计划的支撑必须满足的"不可改进"条件。
    \item \textbf{Kantorovich 对偶}（价格视角）：用"竞争价格对"刻画最小成本，
      给出等价的对偶极大化问题。
  \end{itemize}

  \deflab{本章主结论（定理 5.10）。} 在温和条件下：
  \[
    \min_{\pi\in\coup{\mu}{\nu}}\int_{\X\times\Y} c\dd\pi
    \;=\;
    \sup\Bigl\{\int_{\Y}\ctr{\psi}\dd\nu - \int_{\X}\psi\dd\mu \;:\; \psi\;\text{c-凸}\Bigr\},
  \]
  且极小化器与极大化器均存在，并通过\TLterm{c-次微分}相互刻画。

  \TLtakeaway{本节的核心等式：最优传输成本 $=$ Kantorovich 对偶问题最优值。这是全书第一章的终点，也是后续所有几何应用的出发点。}
\end{frame}

% ── 动机：面包店 / 咖啡馆的循环改进 ─────────────────────────────────────────
\begin{frame}{动机：循环重路由——哪些计划可以改进？（一）}
  \textbf{场景。} 你负责把面包从面包店 $x$ 运到咖啡馆 $y$，已写好一份转运计划。
  有人抱怨成本太高，你决定试试\emph{局部改进}。

  \medskip
  \textbf{改进方法（循环重路由）。}
  \begin{enumerate}\small
    \item 选面包店 $x_1$，它目前向远处咖啡馆 $y_1$ 运货。
    \item 把这一篮面包改道到更近的 $y_2$：节省 $c(x_1,y_2)-c(x_1,y_1)$。
    \item 但 $y_2$ 现在多了面包，再把一篮从面包店 $x_2$ 改道到 $y_3$……
    \item ……直到最终从某个 $x_N$ 把面包改送回 $y_1$，形成一个\emph{闭环}。
  \end{enumerate}

  \medskip
  \deflab{可改进的条件。} 上述重路由使总成本严格降低，当且仅当：
  \[
    c(x_1,y_2)+c(x_2,y_3)+\cdots+c(x_N,y_1)
    \;<\;
    c(x_1,y_1)+c(x_2,y_2)+\cdots+c(x_N,y_N).
  \]

  \TLtakeaway{若支撑中存在可改进的循环，计划\emph{绝对不是最优的}。反之，"无循环可改进"是最优的必要条件——这就是循环单调性。}
\end{frame}

% ── 动机：图示（图 5.1 重绘）─────────────────────────────────────────────────
\begin{frame}{动机：循环重路由图示（图 5.1）}
  \begin{columns}[T]
    \begin{column}{0.55\textwidth}
      \begin{tikzpicture}[scale=0.95,>={Stealth[length=5pt]}]
        % 面包店节点（左侧，圆形）
        \foreach \i/\y in {1/2.4, 2/1.2, 3/0.0}{
          \node[draw=TLprimary,fill=TLprimaryTint,circle,minimum size=16pt,
                inner sep=1pt,font=\scriptsize] (x\i) at (0,\y) {$x_{\i}$};
        }
        % 咖啡馆节点（右侧，矩形）
        \foreach \j/\y in {1/2.4, 2/1.2, 3/0.0}{
          \node[draw=TLaccent,fill=TLaccentLight,rectangle,minimum size=16pt,
                inner sep=2pt,font=\scriptsize] (y\j) at (4.2,\y) {$y_{\j}$};
        }
        % 原始计划（实线）
        \draw[->,TLink,thick] (x1) -- (y1) node[midway,above,font=\scriptsize,TLinkSoft]{原};
        \draw[->,TLink,thick] (x2) -- (y2);
        \draw[->,TLink,thick] (x3) -- (y3);
        % 重路由循环（虚线，橙色）
        \draw[->,TLaccent,dashed,thick] (x1) to[bend left=18] (y2)
              node[near end,above,font=\scriptsize,TLaccent]{改};
        \draw[->,TLaccent,dashed,thick] (x2) to[bend right=18] (y3);
        \draw[->,TLaccent,dashed,thick] (x3) to[bend right=25] (y1);
        % 图例
        \node[font=\scriptsize,TLinkSoft] at (0,-0.6) {面包店};
        \node[font=\scriptsize,TLinkSoft] at (4.2,-0.6) {咖啡馆};
      \end{tikzpicture}
    \end{column}
    \begin{column}{0.42\textwidth}
      \textbf{图例说明：}
      \begin{itemize}\small
        \item 实线 = 原始计划（$x_i\to y_i$）。
        \item 虚线 = 重路由后的循环：$x_1\to y_2,\ x_2\to y_3,\ x_3\to y_1$。
        \item 若虚线总成本 $<$ 实线总成本，原计划可被改进。
        \item 循环单调的计划：对\emph{所有}有限循环，重路由均不能降低成本。
      \end{itemize}
      \vspace{6pt}
      {\scriptsize 据 Villani 图 5.1 重绘。}
    \end{column}
  \end{columns}
  \TLtakeaway{若存在一个循环使重路由更便宜，原计划必然\emph{不最优}。循环单调性精确刻画了"无利可图的循环"。}
\end{frame}

% ── 定义 5.1：c-循环单调性 ────────────────────────────────────────────────────
\begin{frame}{定义 5.1：c-循环单调集与 c-循环单调计划}
  \deflab{定义 5.1（c-循环单调性，Villani Def 5.1）。}
  设 $\X,\Y$ 为任意集合，$c:\X\times\Y\to(-\infty,+\infty]$。
  子集 $\Gamma\subset\X\times\Y$ 称为\TLterm{c-循环单调的}，若对任意 $N\in\N$
  与任意 $(x_1,y_1),\ldots,(x_N,y_N)\in\Gamma$，均有
  \[
    \sum_{i=1}^{N} c(x_i,y_i) \;\le\; \sum_{i=1}^{N} c(x_i,y_{i+1}),
    \qquad (\text{约定 }y_{N+1}=y_1).
  \]
  一个转运计划称为 c-循环单调的，若它集中（concentrated）在某个 c-循环单调集上。

  \begin{itemize}\small
    \item \textbf{直觉：} c-循环单调计划是"无法通过重路由有限次循环来节省成本"的计划——某种意义上的\emph{局部极小化器}。
    \item \textbf{必要性（直觉）：} 最优计划必然 c-循环单调。若不然，存在一个更便宜的循环重路由，与最优性矛盾。
    \item \textbf{充分性（非显然）：} 反过来，c-循环单调是否保证最优？这正是定理 5.10 的核心内容。
  \end{itemize}

  \TLtakeaway{c-循环单调性 = "无法通过循环重排来节省总成本"。它是最优计划的\emph{必要}条件，在温和假设下也是\emph{充分}条件（定理 5.10）。}
\end{frame}

% ── 1D 例子：c-循环单调性与单调重排 ─────────────────────────────────────────
\begin{frame}{可算的例子：1D 中 $c=|x-y|^2$ 的 c-循环单调性}
  \textbf{设定。} $\X=\Y=\R$，$c(x,y)=|x-y|^2$，转运计划 $\pi$。

  \medskip
  \deflab{命题（1D 等价）。} 在此成本下，$\pi$ c-循环单调当且仅当其支撑是\emph{单调不减}的子集（即：$(x_1,y_1),(x_2,y_2)\in\supp\pi$，$x_1<x_2$ 则 $y_1\le y_2$）。

  \medskip
  \textbf{证明简述（$N=2$ 情形）。} 对 $(x_1,y_1),(x_2,y_2)\in\supp\pi$，c-循环单调要求：
  \[
    |x_1-y_1|^2+|x_2-y_2|^2 \;\le\; |x_1-y_2|^2+|x_2-y_1|^2.
  \]
  展开化简得 $(x_1-x_2)(y_1-y_2)\ge0$，即 $x_1,x_2$ 与 $y_1,y_2$ 的大小关系相同——这正是单调性。

  \medskip
  \deflab{推论。} 对 $c=|x-y|^2$，唯一的最优（c-循环单调）计划由\TLterm{单调重排}给出：
  若 $\mu,\nu$ 为 $\R$ 上的概率测度，最优传输映射为 $T=F_{\nu}^{-1}\circ F_{\mu}$（分布函数的反函数复合）。

  \TLtakeaway{1D 二次成本：c-循环单调 $\Leftrightarrow$ 支撑单调不减 $\Rightarrow$ 最优映射 = 单调重排 $T=F_\nu^{-1}\circ F_\mu$。}
\end{frame}

% ── 对偶 Kantorovich 问题：价格视角 ─────────────────────────────────────────
\begin{frame}{价格视角：对偶 Kantorovich 问题（一）}
  \textbf{经济直觉（Villani 的"中介公司"类比）。}
  设一家中介公司愿意接管所有运输，以价格 $\psi(x)$ 在面包店 $x$ 收购面包，
  以价格 $\phi(y)$ 在咖啡馆 $y$ 出售。为保持\emph{竞争性}，必须满足：
  \[
    \phi(y) - \psi(x) \;\le\; c(x,y) \qquad \forall\, x\in\X,\; y\in\Y.
  \]
  中介的利润为 $\int_{\Y}\phi\dd\nu - \int_{\X}\psi\dd\mu$；他们希望\emph{最大化}利润。

  \medskip
  \deflab{对偶 Kantorovich 问题（式 5.3）。}
  \[
    \sup\left\{\int_{\Y}\phi\dd\nu - \int_{\X}\psi\dd\mu \;:\;
    \phi(y)-\psi(x)\le c(x,y),\; \psi\in L^1(\mu),\; \phi\in L^1(\nu)\right\}.
  \]

  \TLtakeaway{"成本视角"（搬多少钱）与"价格视角"（值多少钱）通过对偶相联系：中介能争取的最大利润 $=$ 自运的最小成本。}
\end{frame}

% ── 弱对偶不等式 ────────────────────────────────────────────────────────────
\begin{frame}{价格视角：弱对偶不等式（命题）}
  \deflab{命题（弱对偶，Villani 式 5.4）。}
  对任意竞争价格对 $\phi(y)-\psi(x)\le c(x,y)$ 以及任意耦合 $\pi\in\coup{\mu}{\nu}$，
  \[
    \int_{\Y}\phi\dd\nu - \int_{\X}\psi\dd\mu
    = \int_{\X\times\Y}[\phi(y)-\psi(x)]\dd\pi
    \le \int_{\X\times\Y} c\dd\pi.
  \]
  \textbf{证明。} 等式由边际条件给出（$\int\phi\dd\pi=\int\phi\dd\nu$，$\int\psi\dd\pi=\int\psi\dd\mu$）；
  不等式由约束 $\phi-\psi\le c$ 及 $\pi\ge0$ 直接得到。

  \TLtakeaway{弱对偶：中介能争取的最大利润不超过自运成本。若某对 $(\pi,\psi,\phi)$ 使等号成立，则同时找到了最优计划和最优价格。}
\end{frame}

% ── 弱对偶推论 ───────────────────────────────────────────────────────────────
\begin{frame}{价格视角：弱对偶推论（二）}
  \deflab{推论（弱对偶不等式）。}
  \[
    \sup_{\phi-\psi\le c}\left[\int_{\Y}\phi\dd\nu-\int_{\X}\psi\dd\mu\right]
    \;\le\;
    \inf_{\pi\in\coup{\mu}{\nu}}\int_{\X\times\Y} c\dd\pi.
  \]
  \begin{itemize}
    \item 对偶问题的\emph{任何}可行解给出原始问题的\emph{下界}——这是"弱对偶"。
    \item \textbf{强对偶}（等号成立）是定理 5.10(i) 的深刻内容，需要额外论证。
    \item 实用技巧：若能展示某对 $(\pi,\psi,\phi)$ 使两边都等于同一个值，则两者同时最优。
  \end{itemize}

  \TLtakeaway{弱对偶：$\text{对偶上界} \le \text{原始下界}$。强对偶等号的成立是 Kantorovich 理论的核心定理（定理 5.10）。}
\end{frame}

% ── c-变换 ──────────────────────────────────────────────────────────────────
\begin{frame}{c-变换的定义}
  \textbf{如何找到最优价格？} 给定 $\psi$，最佳的 $\phi(y)$ 应取 $\psi(x)+c(x,y)$ 的下确界：

  \deflab{定义（c-变换）。} 函数 $\psi:\X\to\R\cup\{+\infty\}$ 的\TLterm{c-变换}为
  \[
    \ctr{\psi}(y) \;:=\; \inf_{x\in\X}\bigl[\psi(x)+c(x,y)\bigr],\qquad y\in\Y.
  \]
  （记号 $\ctr{\psi}=\psi^c$。）对称地，$\phi:\Y\to\R\cup\{-\infty\}$ 的 c-变换为
  $\phi^c(x):=\sup_{y\in\Y}[\phi(y)-c(x,y)]$。

  \begin{itemize}
    \item 直觉：$\ctr{\psi}(y)$ 是"在 $y$ 处，考虑所有可能的 $x$ 来源后，诚实竞争价格的上界"。
    \item 由于 $\ctr{\psi}(y)-\psi(x)\le c(x,y)$，$(\psi,\ctr{\psi})$ 自动是竞争价格对。
  \end{itemize}

  \TLtakeaway{c-变换把"买方价格 $\psi$"转化为"最优卖方价格 $\ctr{\psi}$"——是 Legendre 变换的推广，也是对偶理论的核心操作。}
\end{frame}

% ── 紧价格对 ─────────────────────────────────────────────────────────────────
\begin{frame}{紧价格对与对偶问题的约简}
  \deflab{紧价格对（Villani 式 5.5）。} 价格对 $(\psi,\phi)$ 是\TLterm{紧的（tight）}，若
  \[
    \phi(y)=\inf_{x}[\psi(x)+c(x,y)]=\ctr{\psi}(y)
    \quad\text{且}\quad
    \psi(x)=\sup_{y}[\phi(y)-c(x,y)].
  \]
  即：不能再提高卖价，也不能降低买价——价格已达"诚实竞争"的极致。

  \deflab{约简引理。} 从任意竞争价格对 $(\psi,\phi)$ 出发，令 $\phi_1=\ctr{\psi}$，$\psi_1=(\phi_1)^c$，则 $(\psi_1,\phi_1)$ 是紧价格对，且
  $\int\phi_1\dd\nu-\int\psi_1\dd\mu \ge \int\phi\dd\nu-\int\psi\dd\mu$（利润不减）。

  因此，对偶问题可以\emph{不失一般性}地只在紧价格对上取极值。

  \TLtakeaway{一步迭代即可将任意竞争价格对变成紧的：$(\psi,\phi)\mapsto(\psi_1,\ctr{\psi_1})$。对偶问题化为：寻找最优 $\psi$（c-凸性即为其特征）。}
\end{frame}

% ── 定义 5.2：c-凸性 ─────────────────────────────────────────────────────────
\begin{frame}{定义 5.2：c-凸函数}
  \deflab{定义 5.2（c-凸函数，Villani Def 5.2）。}
  设 $c:\X\times\Y\to(-\infty,+\infty]$。函数 $\psi:\X\to\R\cup\{+\infty\}$
  称为\TLterm{c-凸的}，若它不恒为 $+\infty$，且存在 $\zeta:\Y\to\R\cup\{\pm\infty\}$ 使得
  \[
    \psi(x) \;=\; \sup_{y\in\Y}\bigl[\zeta(y)-c(x,y)\bigr] \qquad \forall\,x\in\X.
  \]
  此时 $\psi$ 的 c-变换为 $\ctr{\psi}(y)=\inf_{x}[\psi(x)+c(x,y)]$，
  c-次微分为
  \[
    \partial_c\psi \;:=\; \{(x,y)\in\X\times\Y \;:\; \ctr{\psi}(y)-\psi(x)=c(x,y)\}.
  \]
  $\psi$ 与 $\ctr{\psi}$ 互称为\TLterm{c-共轭函数}。

  \begin{itemize}\small
    \item \textbf{等价描述。} $\psi$ c-凸 $\Leftrightarrow$ $\psi(x)=\inf_{y}[\ctr{\psi}(y)+c(x,y)]$ 对所有 $x$ 成立（即 $\cctr{\psi}=\psi$，见命题 5.8）。
    \item \textbf{c-次微分的含义。} $(x,y)\in\partial_c\psi$ 等价于 $x\in\arg\min_{x'}[\psi(x')+c(x',y)]$——即 $y$ 只接受从\emph{总成本最低}的 $x$ 运来的货物。
  \end{itemize}

  \TLtakeaway{c-凸函数 = "由 $c(x,\cdot)$ 的平移从下方刻画"的函数（图 5.2 直觉）。$\partial_c\psi$ 是最优匹配关系的几何编码。}
\end{frame}

% ── c-凸性图示（图 5.2 重绘）────────────────────────────────────────────────
\begin{frame}{c-凸性的几何直觉（图 5.2 重绘）}
  \begin{columns}[T]
    \begin{column}{0.55\textwidth}
      \begin{tikzpicture}[scale=1.0,>={Stealth[length=4pt]}]
        % x 轴
        \draw[->] (-0.3,0)--(5.2,0) node[right,font=\scriptsize]{$x$};
        \draw[->] (0,-0.3)--(0,3.2) node[above,font=\scriptsize]{值};
        % c-凸函数 psi（折线近似）
        \draw[TLprimary,very thick]
          (0.2,2.4) .. controls (1.2,1.0) and (2.0,0.5) .. (2.6,0.6)
          .. controls (3.2,0.8) and (4.0,1.8) .. (4.8,2.6)
          node[right,font=\scriptsize,TLprimary]{$\psi$};
        % 三个"工具"曲线：zeta(y_i) - c(x, y_i)（细线，从下方贴近）
        \foreach \yi/\col/\pos in {0/TLaccent/0.8, 1/TLpos/2.6, 2/TLinkSoft/4.2}{
          \draw[\col,thin,dashed]
            (0.2,{2.2-0.6*\yi}) .. controls (2,{1.1-0.4*\yi}) and (3,{0.3+0.3*\yi}) .. (4.8,{1.9+0.2*\yi});
        }
        % 接触点（c-次微分元素）
        \fill[TLaccent] (0.9,1.6) circle (2.5pt) node[right,font=\scriptsize,TLaccent]{$y_0\in\partial_c\psi(x_0)$};
        \fill[TLpos] (2.6,0.6) circle (2.5pt) node[right,font=\scriptsize,TLpos]{$y_1\in\partial_c\psi(x_1)$};
        \fill[TLinkSoft] (4.1,1.9) circle (2.5pt) node[left,font=\scriptsize,TLinkSoft]{$y_2\in\partial_c\psi(x_2)$};
      \end{tikzpicture}
    \end{column}
    \begin{column}{0.42\textwidth}
      \textbf{直觉说明：}
      \begin{itemize}\small
        \item 虚线为 $\zeta(y_i)-c(\cdot,y_i)$（形状随 $y_i$ 变化的"工具"）。
        \item $\psi$ 是这些工具的\emph{逐点上确界}。
        \item 每个接触点 $(x_i,y_i)$ 即一个 c-次微分元素：工具恰好从下方"吻到"$\psi$。
        \item Villani："c-凸函数是可以被形状为 $-c(x,\cdot)$ 的工具从下方完整抚摸的函数。"
      \end{itemize}
      \vspace{4pt}
      {\scriptsize 据 Villani 图 5.2 重绘。}
    \end{column}
  \end{columns}
  \TLtakeaway{c-凸函数的图象被一族"成本工具"从下方生成：每个接触点给出一个最优匹配 $(x_i\leftrightarrow y_i)$，即 $\partial_c\psi$ 中的元素。}
\end{frame}

% ── 特殊情形 5.3 与 5.4 ──────────────────────────────────────────────────────
\begin{frame}{c-凸性的两个重要特殊情形（5.3 与 5.4）}
  \deflab{特殊情形 5.3（$c=-x\cdot y$，Villani PC 5.3）。}
  在 $\R^n\times\R^n$ 上取 $c(x,y)=-x\cdot y$，则：
  \begin{itemize}\small
    \item c-变换 $\ctr{\psi}(y)=\inf_{x}[\psi(x)-x\cdot y]=-\sup_{x}[x\cdot y-\psi(x)]$，
      即 $\ctr{\psi}=-\psi^*$（$\psi$ 的\emph{Legendre 变换}的相反数）。
    \item c-凸 $\Leftrightarrow$ $\psi$ 是\emph{凸函数}且下半连续（lsc）。
    \item 这将 Kantorovich 对偶化为经典 Legendre 对偶——后者是前者的特例。
  \end{itemize}

  \medskip
  \deflab{特殊情形 5.4（$c=d$，距离，Villani PC 5.4）。}
  在度量空间 $(\X,d)$ 上取 $c=d$，则：
  \begin{itemize}\small
    \item c-凸 $\Leftrightarrow$ $\psi$ 是\emph{1-Lipschitz 函数}（$|\psi(x)-\psi(x')|\le d(x,x')$）。
    \item 1-Lipschitz 函数的 c-变换是其自身：$\ctr{\psi}=\psi$。
    \item 由此可以得到 Kantorovich–Rubinstein 公式（定理 5.10 的推论）：
      $W_1(\mu,\nu)=\sup_{\psi\text{ 1-Lip}}\int\psi\dd\mu-\int\psi\dd\nu$。
  \end{itemize}

  \TLtakeaway{c-凸性统一了两种经典对偶：$c=-x\cdot y$ 给出 Legendre 变换；$c=d$ 给出 Kantorovich–Rubinstein 公式。}
\end{frame}

% ── 命题 5.8：c-凸性的等价刻画 ────────────────────────────────────────────────
\begin{frame}{命题 5.8：c-凸性的等价刻画与双重 c-变换}
  \deflab{命题 5.8（Villani Prop 5.8）。}
  对任意函数 $\psi:\X\to\R\cup\{+\infty\}$，定义其\TLterm{二次 c-变换}为 $\cctr{\psi}=(\ctr{\psi})^c$，即
  \[
    \cctr{\psi}(x) = \sup_{y\in\Y}\Bigl[\inf_{x'\in\X}\bigl(\psi(x')+c(x',y)\bigr)-c(x,y)\Bigr].
  \]
  则 $\psi$ 是 c-凸的，当且仅当 $\cctr{\psi}=\psi$（即经过两次 c-变换后函数不变）。

  \medskip
  \textbf{证明。} 一般地，对任意 $\phi:\Y\to\R\cup\{-\infty\}$，有 $(\phi^c)^{cc}=\phi^c$。
  这是因为：取 $x'=x$ 可得 $\le$；取 $y'=y$ 可得 $\ge$，两个方向均由定义直接给出。
  由此：若 $\psi=\zeta^c$（c-凸），则 $\cctr{\psi}=\zeta^{ccc}=\zeta^c=\psi$。
  反之若 $\cctr{\psi}=\psi$，则 $\psi=(\ctr{\psi})^c$ 是 $\ctr{\psi}$ 的 c-变换，故 c-凸。

  \medskip
  \deflab{注记 5.9。} 命题 5.8 是 Legendre 对偶的推广：取 $c=-x\cdot y$ 时恢复
  $\psi^{**}=\psi$（凸的充要条件）。

  \TLtakeaway{c-凸 $\Leftrightarrow$ $\cctr{\psi}=\psi$：两次 c-变换后不变。这推广了 Legendre 对偶的"双共轭等于自身"——整个理论的基本代数事实。}
\end{frame}

% ── 定理 5.10（一）：对偶等式陈述 ────────────────────────────────────────────
\begin{frame}{定理 5.10：Kantorovich 对偶（一）——对偶等式}
  \deflab{定理 5.10（Kantorovich 对偶，Villani Thm 5.10）。}
  设 $\X,\Y$ 为 Polish 空间，$\mu\in\Pp(\X)$，$\nu\in\Pp(\Y)$，$c:\X\times\Y\to\R\cup\{+\infty\}$ 下半连续，且
  $c(x,y)\ge a(x)+b(y)$ 对某上半连续的 $a\in L^1(\mu)$，$b\in L^1(\nu)$。则

  \deflab{(i) 存在对偶等式：}
  \[
    \min_{\pi\in\coup{\mu}{\nu}}\int c\dd\pi
    \;=\;\sup_{\substack{\psi\in L^1(\mu),\,\phi\in L^1(\nu)\\\phi-\psi\le c}}\!\!\!\!
      \Bigl[\int_{\Y}\phi\dd\nu-\int_{\X}\psi\dd\mu\Bigr]
    \;=\;\sup_{\psi\text{ c-凸}}\Bigl[\int_{\Y}\ctr{\psi}\dd\nu-\int_{\X}\psi\dd\mu\Bigr].
  \]
  且对偶问题的上确界可只在 $\psi$ 为 c-凸、$\phi=\ctr{\psi}$ 为 c-凹的对上取得。

  \begin{itemize}\small
    \item 注意左端是\emph{最小值}（最优耦合存在，由定理 4.1）。
    \item 右端对偶问题的\emph{上确界}在温和条件下也是最大值（见定理 5.10(iii)）。
  \end{itemize}

  \TLtakeaway{最优传输成本 $=$ Kantorovich 对偶最优值：成本的最小化与价格利润的最大化在最优处\emph{完全吻合}。}
\end{frame}

% ── 定理 5.10（二）：几何刻画 ────────────────────────────────────────────────
\begin{frame}{定理 5.10：Kantorovich 对偶（二）——几何刻画}
  \deflab{(ii) 最优性的等价条件（Villani Thm 5.10(ii)）。}
  设 $c$ 实值，$C(\mu,\nu)<+\infty$。则存在 c-循环单调 Borel 集 $\Gamma\subset\X\times\Y$，
  使得对任意 $\pi\in\coup{\mu}{\nu}$，以下五个命题等价：
  \begin{enumerate}\small
    \item[(a)] $\pi$ 最优（使 $\int c\dd\pi$ 达到最小值）；
    \item[(b)] $\pi$ 是 c-循环单调的（集中在 c-循环单调集上）；
    \item[(c)] 存在 c-凸函数 $\psi$，使得 $\pi$-a.s.\ $\ctr{\psi}(y)-\psi(x)=c(x,y)$；
    \item[(d)] 存在 $\psi:\X\to\R\cup\{+\infty\}$，$\phi:\Y\to\R\cup\{-\infty\}$，满足 $\phi-\psi\le c$ 对全集成立，且等号 $\pi$-a.s.\ 成立；
    \item[(e)] $\pi$ 集中在 $\Gamma$ 上（同一个 $\Gamma$ 对所有最优 $\pi$ 共享）。
  \end{enumerate}

  \TLtakeaway{等价链：$\text{(a)最优} \Leftrightarrow \text{(b)循环单调} \Leftrightarrow \text{(c)集中在}\;\partial_c\psi \Leftrightarrow \text{(d)竞争价格等号}\;\Leftrightarrow\;\text{(e)集中在}\;\Gamma$。这五个刻画是本章最核心的内容。}
\end{frame}

% ── 等价链的实用意义 ─────────────────────────────────────────────────────────
\begin{frame}{定理 5.10 的实用推论（注记 5.13）}
  \textbf{如何用定理 5.10 判断最优性？}

  \begin{itemize}
    \item \textbf{验证 $\pi$ 最优。} 构造竞争价格对 $(\psi,\phi)$（$\phi-\psi\le c$），
      在 $\pi$ 的支撑上验证等号 $\phi(y)-\psi(x)=c(x,y)$。
    \item \textbf{验证 $(\psi,\phi)$ 最优。} 构造支撑上 $\phi-\psi=c$ 的可行耦合 $\pi$。
    \item \textbf{几何。} 最优 $(\psi,\phi)$ 固定后：质量从 $x$ 到 $y$ 当且仅当
      $x\in\arg\min_{x'}[\psi(x')+c(x',y)]$（咖啡馆只接受总成本最低面包店的货）。
  \end{itemize}

  \TLtakeaway{定理 5.10 的实用法则：构造 $(\pi,\psi,\phi)$ 并在支撑上验证 $\phi-\psi=c$，即可同时证明 $\pi$ 和 $(\psi,\phi)$ 均最优。}
\end{frame}

% ── 定理 5.10 等价条件的注意点 ───────────────────────────────────────────────
\begin{frame}{定理 5.10 等价条件辨析}
  % Manual warn-style block (avoids 0.4pt draw-border overflow from \TLwarnblock)
  \begin{tikzpicture}
    \node[draw=TLneg, fill=TLneg!8, rounded corners=3pt,
          inner sep=8pt, text width=\linewidth-17pt, align=justify] {%
      \small\textbf{\color{TLneg}易混淆：私有 c-单调集 vs 共享集 $\Gamma$}\\[0.3em]
      等价条件 (b)：$\pi$ 集中在\emph{某个} c-循环单调集上——该集合随 $\pi$ 变化（私有）。\\
      等价条件 (e)：所有最优 $\pi$ 共享同一集合 $\Gamma=\partial_c\psi$（$\psi$ = 固定的 Kantorovich 势）。
    };
  \end{tikzpicture}

  \begin{itemize}
    \item 若 $c$ 连续，最小共享集 $\Gamma_{\min}$（所有最优支撑的并的闭包）和最大共享集
      $\Gamma_{\max}$（所有最优 $\psi$ 的 c-次微分的交）均有意义。
    \item 注记 5.11：$c$ 不连续时，"集中在 $\Gamma$" 不等同于 "$\mathrm{supp}\,\pi\subset\Gamma$"。
    \item 注记 5.12：$\Gamma$ 一般不唯一；连续成本时可唯一确定 $\Gamma_{\min}$ 和 $\Gamma_{\max}$。
  \end{itemize}

  \TLtakeaway{公共集 $\Gamma$ 编码所有最优计划共享的几何信息；$\Gamma$ 不唯一，任意选取其一均可用于最优性刻画。}
\end{frame}

% ── 证明思路（一）：构造 Kantorovich 势 ──────────────────────────────────────
\begin{frame}{定理 5.10 证明思路（一）：从 c-循环单调集构造势函数}
  \textbf{核心构造（Rüschendorf/Rockafellar）。}
  给定一个 c-循环单调集 $\Gamma$，固定基点 $(x_0,y_0)\in\Gamma$，令 $\psi(x_0)=0$，
  对任意 $x\in\X$ 定义（Villani 式 5.17）：
  \[
    \psi(x) := \sup_{m,\,(x_1,y_1),\ldots,(x_m,y_m)\in\Gamma}
    \bigl[c(x_0,y_0)-c(x_1,y_0)+c(x_1,y_1)-c(x_2,y_1)
    +\cdots+c(x_m,y_m)-c(x,y_m)\bigr].
  \]

  \textbf{关键验证。}
  \begin{itemize}\small
    \item \textbf{$\psi(x_0)=0$。} 由 c-循环单调性，上式对 $(x_1,y_1)=(x_0,y_0)$ 的所有链均 $\le 0$，而取 $m=1,(x_1,y_1)=(x_0,y_0)$ 时为 $0$，故 $\psi(x_0)=0$。
    \item \textbf{$\psi$ 是 c-凸函数。} 通过引入辅助函数 $\zeta(y)=\sup_{\text{链}}\{\cdots+c(x_m,y)\}$，可将 $\psi$ 改写为 $\psi(x)=\sup_{y}[\zeta(y)-c(x,y)]$，这正是 c-凸性定义。
    \item \textbf{$\Gamma\subset\partial_c\psi$。} 对 $(x,y)\in\Gamma$：利用 $\psi$ 的定义，可验证 $\ctr{\psi}(y)=\psi(x)+c(x,y)$。
  \end{itemize}

  \TLtakeaway{构造关键：c-循环单调性保证了 $\psi(x_0)=0$（有界不发散）；由此得到 c-凸函数 $\psi$，其 c-次微分恰好包含 $\Gamma$。}
\end{frame}

% ── 证明思路（二）：对偶等式的建立 ──────────────────────────────────────────
\begin{frame}{定理 5.10 证明思路（二）：对偶等式的完成}
  \textbf{承接前页的构造（$\psi$ c-凸，$\phi=\ctr{\psi}$，$\pi$ 集中在 $\partial_c\psi$）。}

  \textbf{关键计算。} 由边际条件与 $\phi(y)-\psi(x)=c(x,y)$（$\pi$-a.s.）：
  \[
    \int_{\X\times\Y} c(x,y)\dd\pi
    \;=\; \int_{\X\times\Y}[\phi(y)-\psi(x)]\dd\pi
    \;=\; \int_{\Y}\phi\dd\nu - \int_{\X}\psi\dd\mu.
  \]
  因此：$\int c\dd\pi = \int\ctr{\psi}\dd\nu - \int\psi\dd\mu \le \sup_{\psi'\text{ c-凸}}\Bigl[\int\ctr{\psi'}\dd\nu-\int\psi'\dd\mu\Bigr]$，
  但由弱对偶，后者 $\le\int c\dd\pi$。两个方向均成立，故\emph{等号}处处成立。

  \begin{itemize}\small
    \item \textbf{连续成本的情形：} 上述论证通过"从离散逼近 $\mu_n,\nu_n$ 出发，取弱极限"严格化。
    \item \textbf{下半连续成本：} 进一步通过 $c=\lim c_k$（单调逼近，$c_k$ 有界连续）完成。
    \item \textbf{可测性：} $\psi,\ctr{\psi}$ 的可测性需要额外论证（利用 $c$ 的下半连续性和 Luzin 定理），参见 Villani 推论 5.22。本讲义在此处省略该步骤。
  \end{itemize}

  \TLtakeaway{对偶等式的证明方向：先用 c-循环单调集构造 $\psi$，再通过边际条件与等号条件把积分恒等式"锁住"。}
\end{frame}

% ── Rockafellar 定理（经典情形）────────────────────────────────────────────
\begin{frame}{Rockafellar 定理：$c=-x\cdot y$ 的经典情形}
  取 $c(x,y)=-x\cdot y$ 在 $\R^n\times\R^n$ 上，则定理 5.10 化为经典的 Rockafellar 定理：

  \deflab{Rockafellar 定理（经典凸分析）。}
  $\Gamma\subset\R^n\times\R^n$ 是循环单调的（即"凸分析意义下的循环单调"）当且仅当存在凸 lsc 函数 $f:\R^n\to\R\cup\{+\infty\}$，使得 $\Gamma\subset\partial f$（$f$ 的次微分的图象）。

  \textbf{与定理 5.10 的联系：}
  \begin{itemize}\small
    \item "c-循环单调"（$c=-x\cdot y$）$=$ Rockafellar 意义下的"循环单调"。
    \item "c-凸函数"$=$ 凸 lsc 函数（由特殊情形 5.3）。
    \item "c-次微分 $\partial_c\psi$"$=$ 凸分析次微分 $\partial f$（即 $\{(x,y):f(x)+f^*(y)=x\cdot y\}$）。
    \item 定理 5.10 是 Rockafellar 定理的\emph{严格推广}：将凸分析中"凸函数 $\leftrightarrow$ 次微分"替换为"c-凸函数 $\leftrightarrow$ c-次微分"。
  \end{itemize}

  \TLtakeaway{Rockafellar 定理 = 定理 5.10 在 $c=-x\cdot y$ 时的特殊情形：循环单调集 $\Leftrightarrow$ 凸函数的次微分。最优传输把这一经典结果推广到了任意成本函数。}
\end{frame}

% ── 定理 5.10(iii)：最优解存在的条件 ────────────────────────────────────────
\begin{frame}{定理 5.10(iii)：对偶最优解的存在性条件}
  \deflab{(iii) 对偶最优解（Villani Thm 5.10(iii)）。}
  在定理 5.10(i) 的条件下，若还有逐点上界
  \[
    c(x,y) \;\le\; c_{\X}(x)+c_{\Y}(y), \qquad (c_{\X},c_{\Y})\in L^1(\mu)\times L^1(\nu),
  \]
  则对偶问题的\emph{上确界是最大值}（即存在最优价格对），且
  \[
    \min_{\pi\in\coup{\mu}{\nu}}\int c\dd\pi
    \;=\;\max_{\psi\text{ c-凸}}\Bigl[\int_{\Y}\ctr{\psi}\dd\nu-\int_{\X}\psi\dd\mu\Bigr],
  \]
  最优 $\psi$（即 Kantorovich 势）满足 $\psi\in L^1(\mu)$，$\ctr{\psi}\in L^1(\nu)$。

  \textbf{证明思路。} 前页构造的 $\psi$ 在此条件下有界（可积）：利用上界
  $c_{\X}(x)+c_{\Y}(y)$ 建立 $\psi(x)\ge\phi(y_0)-c_{\Y}(y_0)-c_{\X}(x)$（下界为 $\mu$-可积函数），
  $\phi(y)\le\psi(x_0)+c_{\X}(x_0)+c_{\Y}(y)$（上界为 $\nu$-可积函数），因此两者均可积。

  \TLtakeaway{若成本有逐点可积上界，则对偶问题的"sup"变为"max"：Kantorovich 势 $\psi$ 不仅存在，而且可积、最优——三件事同时成立。}
\end{frame}

% ── 稳定性：定理 5.20 ─────────────────────────────────────────────────────────
\begin{frame}{稳定性：最优计划在弱收敛下保持（定理 5.20）}
  \deflab{定理 5.20（最优传输的稳定性，Villani Thm 5.20）。}
  设 $c\in C(\X\times\Y)$ 连续，$\inf c>-\infty$，$(c_k)$ 一致收敛到 $c$，
  $\mu_k\wkto\mu$，$\nu_k\wkto\nu$。对每个 $k$，设 $\pi_k$ 是 $(\mu_k,\nu_k)$ 之间关于 $c_k$ 的最优耦合，
  且 $\int c_k\dd\pi_k<+\infty$。则：
  \begin{enumerate}\small
    \item $(\pi_k)$ 存在弱收敛子列，极限 $\pi\in\coup{\mu}{\nu}$ 是 c-循环单调的。
    \item 若还有 $\liminf_k\int c_k\dd\pi_k<+\infty$，则 $C(\mu,\nu)<+\infty$，$\pi$ 是 $(\mu,\nu)$ 之间的最优耦合。
  \end{enumerate}

  \textbf{证明要点：}
  \begin{itemize}\small
    \item \emph{存在子列：} $\mu_k,\nu_k$ 弱收敛 $\Rightarrow$ 胶着 $\Rightarrow$ Lemma 4.4 给出 $(\pi_k)$ 胶着 $\Rightarrow$ Prokhorov 抽取子列。
    \item \emph{极限 c-循环单调：} 每个 $\pi_k^{\otimes N}$ 集中在 $c_k$-循环单调集上，该集在 $c_k\to c$ 后取极限为 c-循环单调闭集，弱极限继承此性质。
    \item \emph{最优性：} 由 $c\dd\pi\le\liminf c_k\dd\pi_k$ 及定理 5.10(ii) 给出。
  \end{itemize}

  \TLtakeaway{最优传输计划对边际测度的弱收敛稳定：极限仍是最优的。这保证了最优传输是一个"连续"的运算。}
\end{frame}

% ── 推论 5.23：传输映射的稳定性 ──────────────────────────────────────────────
\begin{frame}{推论 5.23：传输映射的稳定性}
  \deflab{推论 5.23（传输映射稳定，Villani Cor 5.23）。}
  设 $T_k:\X\to\Y$ 是关于 $c_k$ 在 $(\mu,\nu_k)$ 之间的最优传输映射（Monge 映射），
  $\nu_k\wkto\nu$。若 $T:\X\to\Y$ 是关于 $c$ 在 $(\mu,\nu)$ 之间的\emph{唯一}最优传输映射，则
  \[
    T_k \;\To\; T \quad \text{依概率}\ (\mu\text{-概率}):
    \quad
    \forall\eps>0,\quad \mu\bigl[\{x:\,d(T_k(x),T(x))\ge\eps\}\bigr]\;\To\;0.
  \]

  \textbf{证明思路。} 由定理 5.20，$\pi_k=(Id,T_k)_{\#}\mu\wkto(Id,T)_{\#}\mu=\pi$。
  利用 Lusin 定理，$T$ 在某紧子集 $K\subset\X$ 上连续（$\mu[K]\ge1-\delta$）；
  再由 $\pi_k\to\pi$ 弱收敛与 $\{d(T_k(x),T(x))\ge\eps\}$ 对应集合的闭性，得到 $\mu$-概率趋向 0。

  \medskip
  \deflab{注记 5.24。} 唯一性假设\emph{不可缺}：若最优映射不唯一，$T_k$ 可能振荡。
  若 $\mu_k\to\mu$ 而不是 $\mu$ 固定，稳定性一般失效（见注记 5.25 中的反例）。

  \TLtakeaway{若极限最优映射\emph{唯一}，则近似问题的最优映射按概率收敛到它——传输的定性性质在弱极限下保持。}
\end{frame}

% ── Brenier 定理（定理 5.30 准则）────────────────────────────────────────────
\begin{frame}{定理 5.30：Monge 问题的可解性准则}
  \deflab{定理 5.30（Monge 问题的可解性准则，Villani Thm 5.30）。}
  在定理 5.10 的条件下，若还满足：
  \begin{enumerate}\small
    \item $C(\mu,\nu)<+\infty$；
    \item 对任意 c-凸函数 $\psi$，使得 $\partial_c\psi(x)$ 包含\emph{多于一个}元素的 $x$ 的集合为 $\mu$-零测。
  \end{enumerate}
  则存在\emph{唯一的}（律意义下）最优耦合，且它是确定性的 Monge 映射：存在可测 $T:\X\to\Y$ 使得 $\pi=(Id,T)_{\#}\mu$ 是唯一的最优耦合。

  \textbf{证明思路。} 存在 c-凸 $\psi$ 使最优 $\pi$ 集中在 $\partial_c\psi$。由条件 (ii)，$\partial_c\psi(x)$ 几乎处处单点，
  可定义 $T(x)$ 为该唯一的 $y$；任何最优 $\pi$ 均集中在 $T$ 的图上。

  \TLtakeaway{定理 5.30 是"c-次微分几乎处处单点 $\Rightarrow$ 最优映射存在且唯一"的精确表述，为 Brenier 定理奠定基础。}
\end{frame}

% ── Brenier 定理的标准应用 ───────────────────────────────────────────────────
\begin{frame}{Brenier 定理：二次成本下的结论}
  \textbf{标准应用。} 取 $c(x,y)=|x-y|^2/2$ 在 $\R^n$ 上，$\mu\in\Pac(\R^n)$（绝对连续）。

  \deflab{Brenier 定理。}
  \begin{itemize}
    \item c-凸函数 = 凸 lsc 函数（由特殊情形 5.3，$c=|x-y|^2/2$ 等价于 $-x\cdot y$）。
    \item 定理 5.30 条件 (ii) 自动满足：$\mu$ 绝对连续时，凸函数的次梯度集 $\mu$-a.e.\ 单点。
    \item 结论：唯一最优 Monge 映射为
      \[
        T = \nabla\varphi, \qquad \varphi\;\text{凸函数},
      \]
      即最优映射是某个凸函数的梯度。
  \end{itemize}

  \TLtakeaway{Brenier 定理：$c=|x-y|^2/2$，$\mu$ 绝对连续 $\Rightarrow$ 最优映射唯一，形如 $T=\nabla\varphi$（凸函数的梯度）。这是最优传输最著名的结果，是 Monge-Ampère 方程的出发点。}
\end{frame}

% ── Brenier：1D 例子 ─────────────────────────────────────────────────────────
\begin{frame}{Brenier 定理：1D 单调重排（具体例子）}
  \textbf{离散例子。} 取 $\mu=\tfrac{1}{3}(\delta_{0}+\delta_{1}+\delta_{2})$，$\nu=\tfrac{1}{3}(\delta_{1}+\delta_{3}+\delta_{5})$，$c=|x-y|^2/2$。

  最优 Monge 映射：$T(0)=1,\;T(1)=3,\;T(2)=5$（单调）。
  总成本 $=\tfrac{1}{3}\cdot\tfrac{1}{2}(1+4+9)=\tfrac{7}{3}$。
  非单调映射如 $T(0)=5,T(1)=3,T(2)=1$ 的成本 $=\tfrac{1}{3}\cdot\tfrac{1}{2}(25+4+1)=5>\tfrac{7}{3}$，更差。

  \textbf{1D 一般情形（$\mu$ 无原子）。} 设 $F_\mu,F_\nu$ 为分布函数，$F_\nu^{-1}$ 为广义逆，则最优映射为
  \[
    T(x) = F_{\nu}^{-1}(F_{\mu}(x)).
  \]
  直觉：把 $\mu$ 的第 $p$ 分位数对应到 $\nu$ 的第 $p$ 分位数——始终单调。

  \TLtakeaway{1D 单调重排 $T=F_\nu^{-1}\circ F_\mu$ 是 Brenier 定理最简单的具体化：最优映射 = 分布函数的广义逆复合，即"按分位数对齐"。}
\end{frame}

% ── Brenier：高斯情形 ─────────────────────────────────────────────────────────
\begin{frame}{Brenier 定理：高斯情形与矩阵几何平均}
  \textbf{高斯情形（$\R^n$）。} 设 $\mu=\mathcal{N}(0,A)$，$\nu=\mathcal{N}(0,B)$，$A,B$ 正定，$c=|x-y|^2/2$。

  最优映射为线性映射 $T(x)=Sx$，其中
  \[
    S = A^{-1/2}\bigl(A^{1/2}BA^{1/2}\bigr)^{1/2}A^{-1/2}.
  \]
  \begin{itemize}\small
    \item $S$ 满足 $S A S = B$（映射把协方差 $A$ 变为 $B$）；
    \item $SA$ 对称正定，故 $S=\nabla\varphi$ 对某凸函数 $\varphi$（确认 Brenier 定理结论）；
    \item 高斯情形的 $W_2$ 距离：$W_2^2(\mu,\nu)=\mathrm{tr}(A+B-2(A^{1/2}BA^{1/2})^{1/2})$；
    \item $S$ 即两协方差矩阵的\emph{几何平均}——高斯情形的 Riemannian 测地线正中点。
  \end{itemize}

  \TLtakeaway{高斯情形：最优映射 $T=Sx$（线性），$S$ 由矩阵几何平均给出。$W_2$ 距离有闭合公式。这是无穷维 Wasserstein 几何在高斯族上的精确计算。}
\end{frame}

% ── 习题 2 · Exercises ────────────────────────────────────────────────────────
\begin{frame}{习题 2 · Exercises（提示见本页，解答见附录）}
  \begin{itemize}
    \item \textbf{习题 2.1}~(\diffI)\ 设 $\X=\Y=\{1,2,3\}$，$c(i,j)=|i-j|$，$\mu=\nu=\frac13\sum_i\delta_i$。
      写出所有耦合中最优的（Monge）耦合 $\pi^\star$，并验证其支撑是 c-循环单调的。
      \begin{itemize}\footnotesize
        \item \hint{最优 Monge 映射为恒等映射；验证对任意 2-元循环 $(i,i),(j,j)$，改道后不减少成本。}
      \end{itemize}
      % SOLUTION: 最优耦合为恒等映射 pi^star = (1/3)(delta_{(1,1)}+delta_{(2,2)}+delta_{(3,3)})，总成本为 0。对任意 (i,i),(j,j) in supp(pi^star)，c(i,i)+c(j,j)=0 <= c(i,j)+c(j,i)=2|i-j|，c-循环单调性成立。

    \item \textbf{习题 2.2}~(\diffII)\ 设 $\psi(x)=\frac12 x^2$ 在 $\R$ 上，$c(x,y)=\frac12|x-y|^2$。
      计算 $\ctr{\psi}(y)$，并确定 $\partial_c\psi$。
      \begin{itemize}\footnotesize
        \item \hint{$\ctr{\psi}(y)=\inf_x[\frac12 x^2+\frac12|x-y|^2]$，对 $x$ 求导令其为零；$\partial_c\psi=\{(x,y):\ctr{\psi}(y)-\psi(x)=c(x,y)\}$。}
      \end{itemize}
      % SOLUTION: ctr{psi}(y) = inf_x [x^2/2 + (x-y)^2/2] = inf_x [x^2 - xy + y^2/2]。导数 = 2x - y = 0，即 x = y/2，代入得 ctr{psi}(y) = y^2/4 - y^2/2 + y^2/2 = y^2/4。验证：ctr{psi}(y) - psi(x) = y^2/4 - x^2/2 = c(x,y) = (x-y)^2/2 iff y^2/4 - x^2/2 = x^2/2 - xy + y^2/2 iff 0 = x^2 - xy + y^2/4 = (x-y/2)^2，即 x=y/2。故 partial_c psi = {(x,y): x = y/2}，即最优映射 T(x)=2x 把 mu 推前到 nu 时对应 y=2x，x=y/2。

    \item \textbf{习题 2.3}~(\diffIII)\ 设 $\mu\in\Pac(\R^n)$，$\nu\in\Pp(\R^n)$，$c=|x-y|^2/2$。
      利用定理 5.30 论证：若两个最优耦合均为 Monge 耦合（对应映射 $T_1,T_2$），则 $T_1=T_2$（$\mu$-a.s.）。
      \begin{itemize}\footnotesize
        \item \hint{两个最优 Monge 耦合均集中在同一个 $\partial_c\psi$（同一个 Kantorovich 势）的图上；利用 $\mu$ 绝对连续时 $\partial_c\psi(x)$ 几乎处处单点。}
      \end{itemize}
      % SOLUTION: 由定理 5.10(ii)，存在 c-凸函数 psi 使得 pi_k = (Id,T_k)_# mu 均集中在 partial_c psi 上，即 T_k(x) in partial_c psi(x) mu-a.s. 对 k=1,2。由 mu 绝对连续以及 c=|x-y|^2/2 的凸性，partial_c psi(x) 对 mu-a.e. x 为单点（这正是定理 5.30 条件(ii)的内容）。故 T_1(x)=T_2(x) mu-a.s.，即最优 Monge 映射在 mu-a.e. 意义下唯一。
  \end{itemize}
\end{frame}

% ── §2 小结 ──────────────────────────────────────────────────────────────────
\begin{frame}{§2 小结：最优传输的三棱柱——循环单调、对偶、势函数}
  本节建立了以下等价链（定理 5.10 的核心）：
  \[
    \pi\text{ 最优}
    \;\Longleftrightarrow\;
    \pi\text{ c-循环单调}
    \;\Longleftrightarrow\;
    \pi\text{ 集中在 }\partial_c\psi
    \;\Longleftrightarrow\;
    \text{竞争价格等号成立}
  \]

  \begin{itemize}\small
    \item \textbf{循环单调性（几何）。} 无法通过重路由有限循环来节省成本。
    \item \textbf{Kantorovich 对偶（价格）。} 最优传输成本 $=$ 最优价格利润；
      最优价格由 c-凸 Kantorovich 势 $\psi$ 给出，$\phi=\ctr{\psi}$。
    \item \textbf{c-次微分（结构）。} 最优计划集中在 $\partial_c\psi$：
      质量从 $x$ 到 $y$ 当且仅当 $\psi(x)+c(x,y)=\ctr{\psi}(y)$。
    \item \textbf{稳定性（定理 5.20）。} 对边际测度弱收敛稳定。
    \item \textbf{Brenier 定理（应用）。} $c=|x-y|^2/2$，$\mu$ 绝对连续 $\Rightarrow$ 唯一最优映射 $T=\nabla\varphi$（凸函数的梯度）。
  \end{itemize}

  \TLtakeaway{§2 是本讲义的"制高点"：三条路通向最优性（几何、价格、势函数），在定理 5.10 处交汇。后续各节（§3--§5）将在此基础上展开度量、曲率和热方程的应用。}
\end{frame}
